\begin{lemma}
\label{lemma-alternating-cech-complex}
In Situation \ref{situation-complex}. Let $M$ be an $A$-module and
denote $\mathcal{F}$ the associated $\mathcal{O}_X$-module. Then
there is a canonical isomorphism of complexes
$$
\Psi : \colim_e \Hom_A(I^\bullet(f_1^e, \ldots, f_r^e), M)
\longrightarrow
\check{\mathcal{C}}_{alt}^\bullet(\mathcal{U}, \mathcal{F})
$$
functorial in $M$ where the differentials on the $\Hom$-complex
are the contragredients of the differentials on
$I^\bullet(f_1^e, \ldots, f_r^e)$.
\end{lemma}

\begin{proof}
Recall that the alternating {\v C}ech complex is the subcomplex
of the usual {\v C}ech complex given by alternating cochains, see
Cohomology, Section \ref{cohomology-section-alternating-cech}.
As usual we view a $p$-cochain in
$\check{\mathcal{C}}_{alt}^\bullet(\mathcal{U}, \mathcal{F})$
as an alternating function $s$ on $\{1, \ldots, r\}^{p + 1}$
whose value $s_{i_0\ldots i_p}$ at $(i_0, \ldots, i_p)$ lies in
$M_{f_{i_0}\ldots f_{i_p}} = \mathcal{F}(U_{i_0\ldots i_p})$.
On the other hand, a $p$-cochain $t$ in
$\Hom^\bullet(I^\bullet(f_1^e, \ldots, f_r^e), M)$
is a map $t : \wedge^{p + 1}(A^{\oplus r}) \to M$.
Write $[i] \in A^{\oplus r}$ for the $i$th basis element and
write
$$
[i_0, \ldots, i_p] = [i_0] \wedge \ldots \wedge [i_p]
\in \wedge^{p + 1}(A^{\oplus r})
$$
For $t$ as above we set
$$
\Psi(t)_{i_0 \ldots i_p} =
(-1)^p \frac{t([i_0, \ldots, i_p])}{f_{i_0}^e\ldots f_{i_p}^e}
$$
It is clear that $\Psi(t)$ is an alternating cochain.
The rule above is compatible with the transition maps
of the system as the transition map
$$
I^\bullet(f_1^e, \ldots, f_r^e) \leftarrow
I^\bullet(f_1^{e + 1}, \ldots, f_r^{e + 1}),
$$
of (\ref{equation-system}) sends $[i_0, \ldots, i_p]$
to $f_{i_0}\ldots f_{i_p}[i_0, \ldots, i_p]$.
It is clear from the description of the localizations
$M_{f_{i_0} \ldots f_{i_p}}$ in
Algebra, Lemma \ref{algebra-lemma-localization-colimit}
that the rule $\Psi$ defines an isomorphism of cochain modules
in degree $p$ in the colimit. To finish the proof we have to show that the map
is compatible with differentials. To see this, for $t$ as above we compute
\begin{align*}
d(\Psi(t))_{i_0 \ldots i_{p + 1}}
& =
\sum\nolimits_{j = 0}^{p + 1} (-1)^j
\Psi(t)_{i_0\ldots \hat i_j \ldots i_{p + 1}} \\
& =
(-1)^p
\sum\nolimits_{j = 0}^{p + 1} (-1)^j t([i_0 \ldots \hat i_j \ldots i_{p + 1}])
(f_{i_0} \ldots \hat f_{i_j} \ldots f_{i_p})^{-e}
\end{align*}
Recall that the differentials on $I^\bullet(f_1^e, \ldots, f_r^e)$
are the negative of the differentials on $K^\bullet(f_1, \ldots, f_r)$.
Thus
\begin{align*}
\Psi(d(t))_{i_0 \ldots i_{p + 1}} & =
(-1)^{p + 1}
d(t)([i_0, \ldots, i_{p + 1}]) (f_{i_0} \ldots f_{i_{p + 1}})^{-e} \\
& =
(-1)^{p + 1}
t(d([i_0, \ldots, i_{p + 1}])) (f_{i_0} \ldots f_{i_{p + 1}})^{-e} \\
& =
(-1)^{p + 1}
t(-\sum\nolimits_{j = 0}^{p + 1}
(-1)^j f_{i_j}^e [i_0, \ldots, \hat i_j, \ldots i_{p + 1}])
(f_{i_0} \ldots f_{i_{p + 1}})^{-e} \\
& =
-(-1)^{p + 1}
\sum\nolimits_{j = 0}^{p + 1}
(-1)^j t([i_0, \ldots, \hat i_j, \ldots i_{p + 1}])
(f_{i_0} \ldots \hat f_{i_j} \ldots f_{i_p})^{-e}
\end{align*}
The two formulas agree concluding the proof.
\end{proof}
